%---------------------------------------------------------------------------%
%->> Frontmatter
%---------------------------------------------------------------------------%
%-
%-> Generate title page
%-
\maketitle% generate Chinese title page
%\MAKETITLE% generate English title page

\intobmk\chapter*{摘要}

作为弥合听障群体与外界沟通的重要桥梁，手语利用手形、动作轨迹、身体姿态和面部表情等多维视觉信息表达丰富的语义内容。基于视觉感知的手语理解技术能够自动识别手语词汇或翻译连续手语表达，在降低听障群体沟通成本、提升公共服务可及性方面具有重要意义。然而，现有手语识别方法通常依赖正面视角或受控环境数据进行训练，而真实应用中的摄像机位置与用户站位往往形成准正面视角，并伴随背景、光照和交互距离变化，导致模型难以泛化到真实场景。围绕上述问题，本研究以准正面鲁棒手语识别为核心，研究如何利用几何先验、语义先验和主动感知机制提升手语理解系统的跨视角泛化能力、环境鲁棒性和自然交互适应性。

针对真实场景中正面到准正面视角迁移带来的性能下降问题，本研究提出几何先验引导的鲁棒手语识别方法，通过双流对比学习在训练阶段引入三维人体结构信息，在不增加推理端计算开销的前提下增强模型对视角扰动和姿态估计噪声的适应能力。针对手语人外观、光照和背景变化带来的环境干扰，本研究进一步提出语义先验引导的高效场景泛化方法，利用生成模型扩充训练样本的外观与场景多样性，并通过关键帧选择和特征传播机制降低生成成本，提升RGB表征在复杂环境下的鲁棒性。面向真实交互应用，本研究构建基于自主移动机器人的自适应手语理解系统，通过主动感知、视角选择、路径规划、手语翻译和实时反馈机制动态优化观测质量，从而进一步提高系统在开放场景中的泛化性和用户体验。已有实验结果表明，基于几何先验的表征学习有助于提升跨视角手语识别性能，语义先验引导的生成式增广能够增强RGB表征在复杂环境下的泛化能力。后续研究将进一步完善移动机器人手语理解系统，在真实场景中验证主动感知、路径规划和交互反馈对翻译质量、系统实时性和用户体验的影响。


\keywords{手语识别；跨视角泛化；移动机器人}
%---------------------------------------------------------------------------%
